\documentclass[11pt,xcolor=svgnames]{beamer}

\usepackage[T1]{fontenc}
\usepackage[spanish]{babel}
\usepackage[utf8]{inputenc}

\usepackage{graphicx}
\usepackage{tikz}

\usepackage{charter}
\usepackage{inconsolata}
\usepackage{helvet}

\usetikzlibrary{arrows,shapes,snakes}

\usepackage{colores}

\usetheme{umbc2}
% \usefonttheme{serif}
\usefonttheme{structurebold}
\usecolortheme[named=Gray]{structure}

\title{oFlute: reconocimiento de señales aplicado al aprendizaje de la flauta
  dulce} 

\author{José Tomás Tocino García}

\institute[Universidad de Cádiz]{Universidad de Cádiz}
% \subject{JOSLUCA 4} 

\date[Sept 2011]{Septiembre de 2011}

\defbeamertemplate*{title page}{customized}[1][]
{
  \vbox{}
  \vfill
  \begin{centering}
    \vspace{0.05cm}
    \includegraphics[width=0.215\textwidth]{imagenes/logotipo} 
    \vspace{0.2cm}
    \begin{beamercolorbox}[sep=8pt,center,#1]{title}
      \usebeamerfont{title}\inserttitle\par%
    \end{beamercolorbox}%
    \vskip0.5em\par
    \begin{beamercolorbox}[sep=0pt,center,#1]{author}
      \usebeamerfont{author}\small\insertauthor\\
      \texttt{<josetomas.tocino@uca.es>}
    \end{beamercolorbox}
    \vspace{0.2cm}
    \begin{beamercolorbox}[sep=0pt,center,#1]{date}
      \usebeamerfont{date}
      {\scriptsize \insertdate }
    \end{beamercolorbox}%\vskip0.5em

    \includegraphics[width=0.12\textwidth]{imagenes/logo_uca} 

    % {\usebeamercolor[fg]{titlegraphic}\inserttitlegraphic\par}
  \end{centering}
  \vfill

}

\setbeamercolor{footline}{fg=white, bg=white}

\begin{document}

% Transparencia de Inicio -> Título
{
  \setbeamertemplate{footline}{}
  \begin{frame}
    \titlepage
  \end{frame}
}
% \logo{\includegraphics[height=2.5cm]{imagenes/logotipo}}
\setbeamertemplate{background}{\includegraphics[width=\paperwidth,height=\paperheight]{imagenes/fondo1.pdf}}
\normalsize

% Transparencia de índice

\begin{frame}{Índice}
  \tableofcontents
\end{frame}


\section{Introducción}

\begin{frame}{Contexto social}
  \begin{center}
    \Large

    Jóvenes en plena simbiosis con las nuevas tecnologías.

    \bigskip
    \pause
    \bigskip

    Las TIC están llegando a los centros educativos.

    \bigskip
    \pause
    \bigskip

    Técnicas docentes basadas en recursos multimedia e informáticos.

  \end{center}
\end{frame}

\begin{frame}{Concepción del proyecto}
  \begin{block}{Idea}
    Hacer un juego educativo.
  \end{block}

  \pause

  \begin{block}{Primera cuestión}
    ¿Sobre qué aspecto educativo? ¿Qué asignatura se beneficia?
  \end{block}

  \vspace{0.5cm}
  \pause

  \begin{center}
    \LARGE Música \\[0.5cm]

    \pause

    \LARGE Aprendizaje de la flauta dulce
  \end{center}
\end{frame}

\begin{frame}{Objetivos}
  \Large
  \begin{itemize}
  \item Adquisición de base de conocimientos. \pause
  \item Módulo de análisis de sonido. \pause
  \item Sistema de interpretación de canciones. \pause
  \item Sistema de lecciones ampliable. \pause
  \item Interfaz de usuario amigable y fluida.
  \end{itemize}
\end{frame}

\begin{frame}{Motivaciones personales}
  \begin{itemize}
  \item Representación digital del sonido.
  \item Programación de audio.\pause

  \item Bases del DSP.
  \item Técnicas básicas de análisis de audio.\pause

  \item Ampliar conocimientos sobre desarrollo de videojuegos.
  \item Aprender nuevas tecnologías.
  \item Aportar al software libre.
  \end{itemize}
\end{frame}

\section{Descripción}
{
\setbeamertemplate{background}{\includegraphics[width=\paperwidth,height=\paperheight]{imagenes/fondo_simple.pdf}}

\begin{frame}{oFlute}
  \begin{columns}
    \column{0.6\textwidth}

    \begin{block}{}
      \begin{center}
        \Large
        Herramienta lúdico-educativa para el aprendizaje de la flauta dulce.

        \medskip
        \medskip

        Interacción del alumno con la flauta en tiempo real.
      \end{center}
    \end{block}

    \column{0.4\textwidth}
    \begin{block}{}
      \begin{center}
        \vspace{-1.5cm}
        \includegraphics[width=0.6\textwidth]{imagenes/logotipo}

        \bigskip
        \bigskip

        \includegraphics[width=\textwidth]{imagenes/imagen_menuPrincipal}
      \end{center}
    \end{block}   
  \end{columns}
\end{frame}


\begin{frame}{Analizador de notas}
  \begin{center}
    \includegraphics[width=0.7\textwidth]{imagenes/imagen_seccionAnalizador}

  Analiza las notas en tiempo real, de forma individual.
  \end{center}
\end{frame}

\begin{frame}{Motor de lecciones}

  \begin{center}
    \includegraphics[width=0.49\textwidth]{imagenes/imagen_seccionLecciones1}\hspace{0.1cm}
    \includegraphics[width=0.49\textwidth]{imagenes/imagen_seccionLecciones2}

    \medskip

    Motor de lecciones con recursos multimedia, totalmente ampliable y personalizable.
  \end{center}

\end{frame}

\begin{frame}{Motor de canciones}

  \begin{center}
    \includegraphics[width=0.49\textwidth]{imagenes/imagen_seccionCanciones1}\hspace{0.1cm}
    \includegraphics[width=0.49\textwidth]{imagenes/imagen_seccionCanciones2}

    \medskip

    Motor de canciones ampliable, permite la interpretación interactiva de canciones.
  \end{center}

\end{frame}
}


\section{Calendario}

\begin{frame}{Planificación}
\textbf{Desarrollo iterativo.} \pause
\begin{enumerate}
\item Adquisición de base de conocimientos. \pause
\item Desarrollo de analizador básico. \pause
\item Interfaz gráfica de usuario. \pause
\item Motor de lecciones. \pause
\item Motor de canciones. \pause
\end{enumerate}
\end{frame}

{
\setbeamertemplate{background}{\includegraphics[width=\paperwidth,height=\paperheight]{imagenes/fondo_simple.pdf}}
\begin{frame}{Diagrama de Gantt}
\includegraphics[angle=270, width=\textwidth]{imagenes/imagen_diagrama_gantt}
\end{frame}
}
\end{document}

